\subsection{Completing the square with the dressed kernel}
\label{q:sec:square}

The whole content of the quintic step is one exact identity: $\widetilde S_{\le5}$ equals $S_{\le5}$ plus a quadratic form in a shifted deviation, which vanishes on shell, plus an explicit $\Sigma_2$-independent term of order $g_s^4$.
This section derives it line by line, then reads off the saddle and shows that it is the resummation of the recursion \eqref{p:eq:hs5:recursion}.
Throughout, $\mc V_4$ is invertible on $\Sym^2\mc H_0/\ker\mc V_4$, and $\mc V_5$ is assumed compatible with the quotient in the sense of the kernel condition \eqref{p:eq:hs5:kernel-condition}; the inverse of the dressed kernel below then exists for small $g_s$ as a Neumann series, and the note assumes it exists at the level of the conformal field theory.

\subsubsection{Collecting the \texorpdfstring{$\Sigma_2$}{Sigma2} terms}
\label{q:sec:square:collect}

The terms of $\caction[\le 4]$ plus \eqref{p:eq:hs5:S5tilde} that contain $\Sigma_2$ are
\begin{equation}
  \label{q:eq:square:F}
  F(\Sigma_2)
    := - \frac{g_s^2}{4!} \, \mc V_4(\Sigma_2, \Sigma_2)
      + \frac{2 g_s^2}{4!} \, \mc V_4(\Sigma_2, \Psi^2)
      + \frac{g_s^3}{5!} \, \mc V_5(\Sigma_2, \Sigma_2, \Psi);
\end{equation}
the kinetic and cubic terms of $\Psi$ are common to $\widetilde S_{\le5}$ and $S_{\le5}$ and drop out of every difference below.
$F$ is a quadratic polynomial in $\Sigma_2$: a quadratic form (the two terms with two $\Sigma_2$) and a linear form (the source).

\subsubsection{Absorbing the coupling into the kernel}
\label{q:sec:square:kernel}

The two quadratic terms combine into a single bilinear form once the coefficients are aligned.
Since $4!/5!=1/5$,
\begin{equation}
  \frac{g_s^3}{5!}
    = \frac{g_s^2}{4!} \cdot \frac{g_s}{5},
\end{equation}
so that
\begin{equation}
  \label{q:eq:square:absorb}
  - \frac{g_s^2}{4!} \, \mc V_4(\Sigma_2, \Sigma_2)
    + \frac{g_s^3}{5!} \, \mc V_5(\Sigma_2, \Sigma_2, \Psi)
    = - \frac{g_s^2}{4!} \, K_\Psi(\Sigma_2, \Sigma_2),
\end{equation}
where
\begin{equation}
  \label{q:eq:square:K}
  \boxed{\;
  \begin{aligned}
    K_\Psi(A, B)
      &
      := \mc V_4(A, B) - \frac{g_s}{5} \, \mc V_5(A, B, \Psi)
      = \braket{A, k_\Psi(B)},
    \\
    k_\Psi
      &
      := p_{2,2} - \frac{g_s}{5} \, p_{3,2}(\Psi, \cdot).
  \end{aligned}
  \;}
\end{equation}
We call $K_\Psi$ the \emph{dressed kernel}: it is the kinetic form of $\Sigma_2$ with the quintic coupling absorbed, so that $\Sigma_2$ propagates in the background of $\Psi$.
It is a small perturbation of $\mc V_4$, symmetric because $\mc V_4$ and $\mc V_5$ are, and linear in $\Psi$ at order $g_s$.
Its inverse on the quotient is
\begin{equation}
  \label{q:eq:square:neumann}
  k_\Psi^{-1}
    = \sum_{n \ge 0} \Big( \frac{g_s}{5} \, p_{2,2}^{-1} \, p_{3,2}(\Psi, \cdot) \Big)^n \, p_{2,2}^{-1},
\end{equation}
a Neumann series which converges for small $g_s$ whenever $p_{2,2}^{-1}$ is bounded: for the order-by-order statements no new invertibility condition arises at this order beyond the one on $p_{2,2}$, while the statements at finite $g_s$ below assume that $K_\Psi$ itself is invertible on the quotient.

\subsubsection{The source through the dressed kernel}
\label{q:sec:square:source}

The source term of \eqref{q:eq:square:F} is still built from the bare $\mc V_4$.
To complete the square with $K_\Psi$ we rewrite it through $K_\Psi$, using \eqref{q:eq:square:K} with $B=\Psi^2$:
\begin{equation}
  \label{q:eq:square:source}
  \mc V_4(\Sigma_2, \Psi^2)
    = K_\Psi(\Sigma_2, \Psi^2) + \frac{g_s}{5} \, \mc V_5(\Sigma_2, \Psi^2, \Psi)
    = K_\Psi(\Sigma_2, \Psi^2) + \frac{g_s}{5} \, v(\Sigma_2),
\end{equation}
where
\begin{equation}
  \label{q:eq:square:v}
  v(A)
    := \mc V_5(A, \Psi^2, \Psi)
    = \braket{A, p_{3,2}(\Psi^3)}
\end{equation}
is a linear form on $\Sym^2\mc H_0$, cubic in $\Psi$ (two powers through $\Psi^2$, one explicit).
Inserting \eqref{q:eq:square:absorb} and \eqref{q:eq:square:source} into \eqref{q:eq:square:F},
\begin{equation}
  \label{q:eq:square:F-dressed}
  F(\Sigma_2)
    = - \frac{g_s^2}{4!} \Big[
        K_\Psi(\Sigma_2, \Sigma_2)
        - 2 K_\Psi(\Sigma_2, \Psi^2)
        - \frac{2 g_s}{5} \, v(\Sigma_2)
      \Big].
\end{equation}

\subsubsection{The shifted saddle and the square}
\label{q:sec:square:saddle}

The two linear terms in the bracket of \eqref{q:eq:square:F-dressed} are both of the form $2K_\Psi(\Sigma_2,\cdot)$, since $v(\Sigma_2)=K_\Psi(\Sigma_2,K_\Psi^{-1}v)$ by definition of $K_\Psi^{-1}v$:
\begin{equation}
  2 K_\Psi(\Sigma_2, \Psi^2) + \frac{2 g_s}{5} \, v(\Sigma_2)
    = 2 K_\Psi(\Sigma_2, \Sigma_2^*),
\end{equation}
where the shifted saddle is
\begin{equation}
  \label{q:eq:square:saddle}
  \boxed{\;
  \Sigma_2^*
    := \Psi^2 + \frac{g_s}{5} \, K_\Psi^{-1} v
    = \Psi^2 + \frac{g_s}{5} \, k_\Psi^{-1} \, p_{3,2}(\Psi^3).
  \;}
\end{equation}
Let's complete the square.
With the shifted deviation
\begin{equation}
  \label{q:eq:square:Xtilde}
  \tilde X
    := \Sigma_2 - \Sigma_2^*
    = X - \frac{g_s}{5} \, K_\Psi^{-1} v,
\end{equation}
the symmetry of $K_\Psi$ gives
\begin{equation}
  \label{q:eq:square:square}
  K_\Psi(\Sigma_2, \Sigma_2) - 2 K_\Psi(\Sigma_2, \Sigma_2^*)
    = K_\Psi(\tilde X, \tilde X) - K_\Psi(\Sigma_2^*, \Sigma_2^*).
\end{equation}
The variation of $F$ is now $\delta F=-\frac{2g_s^2}{4!}K_\Psi(\delta\Sigma_2,\tilde X)$, so the equation of motion of $\Sigma_2$ is exactly
\begin{equation}
  \tilde X = 0,
  \qquad\text{i.e.}\qquad
  \Sigma_2 = \Sigma_2^* .
\end{equation}

\subsubsection{Reassembling the effective vertices}
\label{q:sec:square:reassemble}

The constant term $K_\Psi(\Sigma_2^*,\Sigma_2^*)$ contains the target.
Expanding \eqref{q:eq:square:saddle} and using $K_\Psi(\Psi^2,K_\Psi^{-1}v)=v(\Psi^2)$ and $K_\Psi(K_\Psi^{-1}v,K_\Psi^{-1}v)=v(K_\Psi^{-1}v)$,
\begin{equation}
  \label{q:eq:square:const}
  K_\Psi(\Sigma_2^*, \Sigma_2^*)
    = K_\Psi(\Psi^2, \Psi^2)
      + \frac{2 g_s}{5} \, v(\Psi^2)
      + \frac{g_s^2}{25} \, v(K_\Psi^{-1} v).
\end{equation}
The first two terms reassemble the quartic and quintic vertices with their coefficients.
From \eqref{q:eq:square:K} and \eqref{q:eq:square:v}, $K_\Psi(\Psi^2,\Psi^2)=\mc V_4(\Psi^4)-\frac{g_s}{5}\mc V_5(\Psi^5)$ and $v(\Psi^2)=\mc V_5(\Psi^5)$, hence
\begin{equation}
  K_\Psi(\Psi^2, \Psi^2) + \frac{2 g_s}{5} \, v(\Psi^2)
    = \mc V_4(\Psi^4) + \frac{g_s}{5} \, \mc V_5(\Psi^5),
\end{equation}
and multiplying by the prefactor $\frac{g_s^2}{4!}$,
\begin{equation}
  \frac{g_s^2}{4!} \, \mc V_4(\Psi^4) + \frac{g_s^2}{4!} \cdot \frac{g_s}{5} \, \mc V_5(\Psi^5)
    = \frac{g_s^2}{4!} \, \mc V_4(\Psi^4) + \frac{g_s^3}{5!} \, \mc V_5(\Psi^5),
\end{equation}
which are exactly the quartic and quintic terms of $S_{\le5}$.
This is how the normalization $1/5!$ of \eqref{p:eq:hs5:S5tilde} is forced: the quintic coefficient of the target is produced by the cross term $\frac{2g_s}{5}v(\Psi^2)$ of the square minus the $-\frac{g_s}{5}\mc V_5(\Psi^5)$ inside $K_\Psi(\Psi^2,\Psi^2)$, and the two combine to $+\frac{g_s}{5}$ only if the coupling carries $1/5!$.
The last term of \eqref{q:eq:square:const} carries the prefactor
\begin{equation}
  \frac{g_s^2}{4!} \cdot \frac{g_s^2}{25}
    = \frac{g_s^4}{600},
  \qquad
  600 = 4! \cdot 25,
\end{equation}
the square of the shift $\frac{g_s}{5}$ against the $\frac{g_s^2}{4!}$ of the kinetic term.

\subsubsection{The identity}
\label{q:sec:square:identity}

Putting \eqref{q:eq:square:F-dressed}, \eqref{q:eq:square:square} and \eqref{q:eq:square:const} together, and adding back the common $\Psi$ terms,
\begin{equation}
  \label{q:eq:square:identity}
  \boxed{\;
  \widetilde S_{\le 5}
    = S_{\le 5}
      - \frac{g_s^2}{4!} \, K_\Psi(\tilde X, \tilde X)
      + \frac{g_s^4}{600} \, v(K_\Psi^{-1} v).
  \;}
\end{equation}
This is an identity in the fields $(\Psi,\Sigma_2)$, valid at every order in $g_s$; it has been checked symbolically as such.
The three terms are the target, a quadratic form in the shifted deviation, and a term independent of $\Sigma_2$.
We call the last one the \emph{feedback}: it is the surplus that the Gaussian elimination of $\Sigma_2$ produces beyond the vertex it was introduced to reproduce.
It is of order $g_s^4$ and of degree six in $\Psi$ ($v$ is cubic and appears twice), so it is a six-point interaction; it does not affect the matching at order $g_s^3$, and it is the exact form of the statement that the auxiliary field feeds back into the action only one order later.

On shell, $\tilde X=0$ and
\begin{equation}
  \label{q:eq:square:onshell}
  \widetilde S_{\le 5}\big(\Psi, \Sigma_2^*\big)
    = S_{\le 5}(\Psi) + \frac{g_s^4}{600} \, v(K_\Psi^{-1} v),
\end{equation}
which is \eqref{p:eq:hs5:eliminate} with the remainder made explicit.
All tree amplitudes of $\widetilde S_{\le5}$ therefore agree with those of $S_{\le5}$ through order $g_s^3$; \cref{p:fig:hs5:five-point} shows the mechanism diagram by diagram.

\subsubsection{The saddle and the recursion}
\label{q:sec:square:recursion}

The saddle \eqref{q:eq:square:saddle} is the resummed form of the recursive solution \eqref{p:eq:hs5:recursion}.
The equation of motion \eqref{p:eq:hs5:eom-sigma}, up to an overall sign $0=p_{2,2}(\Sigma_2)-p_{2,2}(\Psi^2)-\frac{g_s}{5}p_{3,2}(\Psi,\Sigma_2)$, is $k_\Psi(\Sigma_2)=p_{2,2}(\Psi^2)$ in the notation \eqref{q:eq:square:K}, so that
\begin{equation}
  \label{q:eq:square:saddle-R}
  \Sigma_2^*
    = k_\Psi^{-1} \, p_{2,2}(\Psi^2).
\end{equation}
To see that this agrees with \eqref{q:eq:square:saddle}, apply \eqref{q:eq:square:K} to $\Psi^2$: $p_{2,2}(\Psi^2)=k_\Psi(\Psi^2)+\frac{g_s}{5}p_{3,2}(\Psi,\Psi^2)$, and $p_{3,2}(\Psi,\Psi^2)=p_{3,2}(\Psi^3)$ by symmetry; applying $k_\Psi^{-1}$ gives $\Sigma_2^*=\Psi^2+\frac{g_s}{5}k_\Psi^{-1}p_{3,2}(\Psi^3)$.
Expanding $k_\Psi^{-1}$ with the Neumann series \eqref{q:eq:square:neumann},
\begin{equation}
  \label{q:eq:square:recursion}
  \Sigma_2^*
    = \Psi^2
      + \frac{g_s}{5} \, p_{2,2}^{-1} \, p_{3,2}(\Psi^3)
      + \frac{g_s^2}{25} \, p_{2,2}^{-1} \, p_{3,2}\Big( \Psi, p_{2,2}^{-1} \, p_{3,2}(\Psi^3) \Big)
      + O(g_s^3),
\end{equation}
which is the recursion \eqref{p:eq:hs5:recursion} iterated once: each iteration inserts one more $p_{3,2}(\Psi,\cdot)$ between two $p_{2,2}^{-1}$, with one more factor $\frac{g_s}{5}$.
The dressed-kernel inverse is thus the closed form of the recursion: what looks like interactions at all orders generated by $\Sigma_2$ is the geometric series of the self-energy insertions $p_{3,2}(\Psi,\cdot)$ on the $\Sigma_2$ line.
The same series, read in the action rather than in the saddle, expands the feedback as
\begin{equation}
  \frac{g_s^4}{600} \, v(K_\Psi^{-1} v)
    = \frac{g_s^4}{600} \Big[
      v(\mc V_4^{-1} v)
      + \frac{g_s}{5} \, v\big(\mc V_4^{-1} W_\Psi \mc V_4^{-1} v\big)
      + O(g_s^2)
    \Big],
\end{equation}
with $W_\Psi:=\mc V_5(\cdot,\cdot,\Psi)$, where $\mc V_4^{-1}W_\Psi\mc V_4^{-1}v$ is shorthand for $p_{2,2}^{-1}p_{3,2}(\Psi,p_{2,2}^{-1}p_{3,2}(\Psi^3))$.

